\newpage
%\addcontentsline{toc}{chapter}{\MakeUppercase{Implementation}}
\thispagestyle{empty}
%
\begin{spacing}{1.2}
%
\chapter{Implementation}

This chapter explains how the GramaGIS platform was built and integrated, covering data preparation, database configuration, service publishing, backend routes, natural language processing, and the browser-based user interface. Together, these implementation steps realize the system's goal of delivering a practical and maintainable village governance platform.

\section{Data Acquisition and Preparation}

The foundational step for GramaGIS involved the digitization of village-level infrastructure. Using QGIS, spatial features such as ward boundaries, public health centers, educational institutions, road networks, government offices, and amenity points were created as vector layers (polygons, points, and lines) \cite{smart_panchayat2023}.

Each feature was populated with administrative attribute data, including building names, ward numbers, facility types, and operational details.

\subsection*{Coordinate Reference System and Validation}

All datasets were standardized using the \textbf{WGS 84} coordinate reference system (\textbf{EPSG:4326}). This ensures native compatibility with web-mapping standards and GeoServer publishing requirements \cite{fu2020webgis}. After topology validation and attribute verification, the datasets were migrated into the PostgreSQL environment.

\section{Spatial Database Configuration}

The backend storage is powered by \textbf{PostgreSQL} with the \textbf{PostGIS} spatial extension enabled. This allows the database to treat geographic shapes as a native data type \cite{postgis_in_action}.

\subsection*{Spatial Table Design and Indexing}

\begin{itemize}
    \item \textbf{Geometry Columns:} Defined using appropriate geometry types such as \texttt{POINT}, \texttt{LINESTRING}, and \texttt{POLYGON}.
    \item \textbf{Spatial Indexing:} \textbf{GiST (Generalized Search Tree)} indexes were implemented on geometry columns to improve spatial query performance \cite{postgis_in_action}.
    \item \textbf{Application Tables:} Additional relational tables were used for login data and feedback management.
\end{itemize}

\section{GeoServer Service Publishing}

GeoServer acts as the map engine, converting PostGIS data into web-renderable layers and editable services.

\subsection*{Configuration Workflow}

\begin{itemize}
    \item \textbf{Workspace \& Store:} A dedicated \texttt{gramagis} workspace was created and linked to the PostGIS data store.
    \item \textbf{Layer Styling:} Layer styling and naming conventions were configured to provide clear symbology for village assets.
    \item \textbf{Service Exposure:} Layers were published as \textbf{WMS} for visual rendering, \textbf{WFS} for feature retrieval, and \textbf{WFS-T} for controlled editing \cite{fu2020webgis}.
\end{itemize}

\section{Backend API and Proxy Development}

The application logic was developed using \textbf{Node.js} and the \textbf{Express.js} framework. This layer acts as the main application backend between the client, database, GeoServer workflows, and the AI service \cite{fu2020webgis}.

\subsection*{Express.js Routing and Proxying}

\begin{itemize}
    \item \textbf{Protected Proxy Routes:} Express.js routes were created for WFS access, schema inspection, dataset download, and WFS-T edit operations. These routes append internal GeoServer credentials when required and shield administrative operations from public access.
    \item \textbf{JWT Authentication:} A middleware layer using JSON Web Tokens (JWT) was implemented to verify the identity of Panchayat officials before allowing access to administrative features \cite{smart_panchayat2023}.
    \item \textbf{Database-Backed Login:} Authentication credentials are validated against a login table in PostgreSQL, with password hashes verified on the backend before issuing a JWT.
\end{itemize}

\section{Natural Language Query Implementation}

GramaGIS includes an NL2CQL (Natural Language to Command Query Language) translation module implemented within the Express.js environment \cite{liu2025nalspatial}.

\subsection*{Prompt Construction and Translation}

When a user submits a search string, the frontend and backend coordinate to provide the Gemini API with a schema-aware prompt. The prompt contains the relevant layer names and attribute fields (e.g., \texttt{ward\_no}, \texttt{ownership}, \texttt{type}). The API translates the natural language input into a structured JSON response containing a target layer and a \textbf{CQL\_FILTER} string \cite{liu2025nalspatial}.

\subsection*{Query Execution Flow}

\begin{itemize}
    \item \textbf{Input Processing:} The user input is captured by the frontend and sent to the Express.js server.
    \item \textbf{Filter Generation:} The Gemini API returns a JSON object containing the target layer and the CQL filter logic \cite{liu2025nalspatial}.
    \item \textbf{Frontend Filter Application:} The frontend applies this CQL string to the active map layer and also issues a filtered WFS request through the Express proxy.
    \item \textbf{Spatial Filtering:} GeoServer applies the filter to the underlying spatial services, returning only the geometries and attributes that match the user's criteria.
\end{itemize}

\section{Frontend Development with Leaflet}

The user interface was built as a web application using \textbf{Leaflet.js} for spatial rendering.

\subsection*{Map Interaction and Administrative Tools}

\begin{itemize}
    \item \textbf{Layer Control:} A custom UI allows users to toggle village layers dynamically.
    \item \textbf{Transactional Editing:} For administrators, a custom mini-map geometry editor was implemented for drawable layers. Modified features are sent to the backend as attribute data and WKT geometry, which the server converts into \textbf{WFS-T} XML requests to update the spatial database \cite{fu2020webgis, postgis_in_action}.
    \item \textbf{Metadata Retrieval:} Feature-click listeners were implemented to trigger live \texttt{GetFeatureInfo} requests for asset-specific details, which are then displayed in a formatted panel.
\end{itemize}

\end{spacing}
